\documentclass[12pt]{article}
\usepackage[utf8]{inputenc}
\usepackage[T1]{fontenc}
\usepackage[margin=1in]{geometry}
\usepackage{setspace}
\usepackage{xcolor}
\usepackage[colorlinks=true, urlcolor=blue!60!black, linkcolor=blue!60!black]{hyperref}
\setlength{\parindent}{0em}
\setlength{\parskip}{0.45em}
\renewcommand{\baselinestretch}{1.0}
\begin{document}
\begin{center}
{\Large \textit{Calder v. Bull}, 3 U.S. (3 Dall.) 386 (1798)\par}
\vspace{0.3em}
{\normalsize Case Brief --- One-Page Edition\par}
\end{center}

\textbf{Memory Jogger.} The Ex Post Facto Clause bars only retroactive criminal laws (\href{https://supreme.justia.com/cases/federal/us/3/386/}{Justia}).

\textbf{Facts.} After the Bulls' time to appeal a probate decree disapproving Morrison's will lapsed, Connecticut's legislature, exercising its customary judicial power, granted a new hearing, and the will was approved. The heirs-at-law, the Calders, invoked Article I, \S~10.

\textbf{Procedural History.} Connecticut's courts upheld the resolution; writ of error to the Supreme Court.

\textbf{Votes.} Affirmed 4--0 in \textit{seriatim} opinions (Chase, Paterson, Iredell, Cushing).

\textbf{Holding.} The Ex Post Facto Clauses reach only retroactive criminal laws in Chase's four categories; the civil resolution was valid.

\textbf{Reasoning.} ``Ex post facto'' was a technical term for penal laws, and reading it to cover civil laws would make the Contracts and Takings Clauses superfluous. Chase added that natural-law principles limit legislatures even without text; Iredell answered that courts may void only laws violating the written Constitution. Both views were dicta.

\textbf{Examples.} \emph{Same-Side:} a retroactive civil tax fix. \emph{Opposite-Side:} retroactively dropping corroboration (\textit{Carmell}). \emph{Fence-Sitter:} sex-offender registration (\textit{Smith v. Doe}).

\textbf{Critique.} Chase condemned taking from A to give to B yet upheld a shifted inheritance; Justice Thomas questions the criminal-only limit.

\textbf{Key Quotation.} ``1st. Every law that makes an action done before the passing of the law and which was innocent when done, criminal and punishes such action. 2d. Every law that aggravates a crime, or makes it greater than it was when committed. 3d. Every law that changes the punishment, and inflicts a greater punishment than the law annexed to the crime, when committed. 4th. Every law that alters the legal rules of evidence and receives less or different testimony than the law required at the time of the commission of the offence in order to convict the offender.'' --- Justice Chase, 3 U.S. at 390.

\end{document}
